% Generated TeX body; no PDF or document-style claim.
% Maintained source: problems/general-matching/reference/certificate-proof.md
% Recovered-source SHA-256: fe6772bb46d8cfb889d89acc0556072d44fe845417e5b5c72c6fc6c2015468d5
% Rebuild: python3 tools/build_manuscripts.py
% Common packages: amsmath, amssymb, hyperref, longtable, booktabs, array.
% Shared original Subset Sum preamble: manuscripts/paper-style.tex.

\hypertarget{af-43f6530146-a-quadratic-extraction-of-a-matching-optimality-certificate}{%
\section{A quadratic extraction of a matching optimality certificate}\label{af-43f6530146-a-quadratic-extraction-of-a-matching-optimality-certificate}}

This note proves the certificate extraction implemented by \texttt{\_failed\_search\_forest} and \texttt{\_barrier} in \href{https://github.com/icekylinx/algorithm-frontier/blob/main/problems/general-matching/reference/matching.py}{matching.py}. Given a maximum matching of size \(\nu\) on a simple graph with \(k\) vertices and \(h\) edges, extraction takes \(O(h+k(\nu+1))\) additional work, at most \(O(k^2)\). This preserves the reference solver's cubic bound on general components. The certificate is checked by the \href{https://github.com/icekylinx/algorithm-frontier/blob/main/problems/general-matching/reference/verify_certificate.py}{independent verifier}, which performs no blossom search.

The public solver uses this extraction on nonbipartite retained kernels. On a bipartite kernel, a \href{https://github.com/icekylinx/algorithm-frontier/blob/main/problems/general-matching/reference/bipartite-proof.md\#5-a-linear-certificate-on-the-first-kernel}{separate local proof} constructs a minimum vertex cover in \texttt{O(k+h)} work and proves that it is a tight certificate in the same format. This avoids a cubic certificate cost in the faster bipartite branch. The theorem below continues to apply to \texttt{\_failed\_search\_forest} and \texttt{\_barrier} on any simple graph.

The construction is a classical alternating-forest method. The resulting sets give the Gallai--Edmonds decomposition, whose historical source is \href{https://real.mtak.hu/189486/1/cut_MATKUTINT_1964_3_pp401_-_413.pdf}{Gallai (1964), Theorem (E.1)}. The proof below establishes the required extraction and upper-bound equality directly; it does not assume that decomposition theorem.

\hypertarget{af-43f6530146-statement-and-algorithm}{%
\subsection{1. Statement and algorithm}\label{af-43f6530146-statement-and-algorithm}}

Fix a maximum matching \(M\) of the simple graph \(H\). Let \(f\) be the number of vertices exposed by \(M\), meaning unmatched by it. Start one Edmonds alternating forest with all exposed vertices as roots. Keep \(M\) fixed throughout the search.

Every root is initially \textbf{outer}. An edge from an outer vertex to an unreached matched vertex grows the forest by that \textbf{inner} vertex and its matched outer child. An edge between different outer blocks in the same tree contracts the usual odd blossom. Every original vertex in the new block becomes outer; newly outer vertices enter the scan queue. An edge between outer vertices in different trees would give an augmenting path, so the implementation rejects that state as inconsistent with a maximum input matching.

After exhausting the queue, let \(O\) be the set of all original outer vertices and \(I\) the remaining inner vertices. Return \(A=I\).

\textbf{Theorem 1.} The algorithm leaves \(M\) unchanged and returns a set \(A\) satisfying

\[
2|M|=|V(H)|-o(H-A)+|A|,
\]

where \(o(H-A)\) is the number of odd connected components after deleting \(A\). It uses \(O(h+k(\nu+1))=O(k^2)\) time and \(O(k+|E(H)|)\) storage. Moreover,

\[
O=\{v:\nu(H-v)=\nu(H)\},
\qquad I=N_H(O)\setminus O.
\]

The implementation first tests \(A=\varnothing\) and skips the forest search when that set already gives the displayed equality. This shortcut returns a valid barrier, although it need not be the particular set \(N_H(O)\setminus O\). The theorem describes the forest search when it is used; the independent verifier accepts either tight barrier.

\hypertarget{af-43f6530146-invariants-of-the-contracted-forest}{%
\subsection{2. Invariants of the contracted forest}\label{af-43f6530146-invariants-of-the-contracted-forest}}

A \textbf{block} is either a singleton or the set of original vertices inside a current contracted blossom. The quotient graph has one vertex per block. At every stage the search maintains these invariants:

\begin{enumerate}
\def\labelenumi{\arabic{enumi}.}
\tightlist
\item
  Each outer block is connected, has odd order, and is factor-critical: deleting any prescribed vertex leaves a perfect matching inside that block. Its current \(M\)-edges form a near-perfect internal matching whose only internally exposed vertex is the block's base.
\item
  \textbf{Only the base can have a matched edge leaving an outer block.} A nonroot outer block's base is matched to its inner parent. A root block has no external matched edge and contains exactly one original exposed vertex of \(M\).
\item
  Every inner vertex is a singleton and has exactly one matched outer child. Thus every nonroot outer block has one inner parent, and every inner vertex has one such child. There is one more outer block than inner vertices in each quotient tree.
\item
  The original members of different blocks are disjoint. Every original member of an outer block is marked outer and has its adjacency scanned once after receiving that mark.
\end{enumerate}

These statements initially hold for the singleton roots. Growing a tree adds one inner singleton and its matched outer child, preserving the matching structure and the difference between the two node counts.

For a contraction, take an edge between outer blocks in the same tree. The least common tree ancestor of its endpoints is outer: an inner vertex has only one child and cannot be the branching ancestor. The two paths to this ancestor, together with the new edge, form an odd cycle in the quotient, with one more outer block than inner singleton. If the cycle contains \(a\) outer blocks, it contains \(a-1\) inner vertices. Replacing them by one outer block removes \(a-1\) of each type, preserving the count exactly.

The new block is factor-critical. More generally, consider an odd cycle of disjoint odd factor-critical blocks joined by selected connecting edges. After deleting a vertex in one block, match that punctured block internally. Pair the remaining even number of blocks around the cycle. For each pair use its connecting edge, and match both blocks internally after deleting the connecting endpoints. This is a perfect matching of the union minus the prescribed vertex. The argument applies recursively to nested blossoms and uses only actual graph edges.

The fixed matching also has the required base property after contraction. Each inner vertex on the alternating cycle is matched to its outer child on that cycle. Each constituent outer block is internally matched except at its base. Consequently all bases on the cycle are matched internally to the new block except the ancestor block's base. That base retains its previous external matched edge, if any. No second external matched edge appears.

For the array implementation, \texttt{base{[}v{]}} identifies the current quotient block. To climb from an outer block to its outer grandparent, first use the matched edge at its base and then the \texttt{parent} edge of that inner vertex. These are exactly the two steps used in \texttt{common\_base}. \texttt{mark\_branch} rewrites the alternating routes along the two cycle branches; the later scan identifies all their blocks with the ancestor base. The ancestor's external matched edge is unchanged. Tree branches attached to the contracted cycle remain attached to its new block. The root label \texttt{owner} remains unchanged within that tree. This maintains the quotient forest and the unique external matched-edge invariant through every parent rewrite.

\hypertarget{af-43f6530146-lifting-paths-through-the-blocks}{%
\subsection{3. Lifting paths through the blocks}\label{af-43f6530146-lifting-paths-through-the-blocks}}

Every vertex of an outer block can be exposed by a matching of the same cardinality as \(M\), preserving the other roots. To see this, follow the quotient tree path from its root block to the target block. Remove the matched edge from each inner vertex to the next block base, and add the preceding unmatched tree edge into that inner vertex.

In each nonfinal block on this path, the new crossing edge uses one known attachment vertex \(x\). Replace the internal matching by a perfect matching of that block minus \(x\), using factor-criticality. In the final block leave the prescribed target vertex exposed. Other blocks and off-path matching edges stay as they were. Exactly one original root has been matched and the target has been exposed; cardinality is unchanged.

The symmetric difference with \(M\) contains a simple even alternating path from that root to the target. Its first edge is unmatched by \(M\) and its last edge is matched by \(M\); a zero-length path is allowed. This is the required lifting argument. It does not assume that an arbitrary projected walk through nested blossoms is already simple.

If an edge joined outer vertices in different trees, the two lifted paths would be vertex-disjoint because their trees have disjoint blocks. Joining them by that edge gives an augmenting path between two exposed roots. Such a path would increase \(|M|\), contradicting maximal cardinality. This proves that the explicit cross-tree rejection cannot occur for a correct maximum matching.

\hypertarget{af-43f6530146-exhaustion-gives-a-tight-odd-component-bound}{%
\subsection{4. Exhaustion gives a tight odd-component bound}\label{af-43f6530146-exhaustion-gives-a-tight-odd-component-bound}}

Let \(b\) be the number of final outer blocks. Summing the quotient-tree count over all \(f\) roots gives

\[
b=|I|+f.
\]

There is no edge between distinct final outer blocks. Such an edge would have been scanned when its later-labelled endpoint became outer. It would contract blocks within a tree or give an augmenting path between trees. Later contractions only merge blocks.

Also \(N_H(O)\setminus O=I\). An unreached neighbor would grow a tree, while every inner vertex has an outer parent. Let \(U=V(H)\setminus(O\cup I)\). Every vertex in \(U\) is matched by \(M\) to another vertex of \(U\): all exposed vertices were roots, and every reached matched pair lies inside an outer block or is an inner/outer parent-child pair. Hence \(M\) restricted to \(U\) is a perfect matching.

After deleting \(I\), the components are the \(b\) connected odd outer blocks and the components of \(H[U]\). Each latter component has a perfect matching and therefore even order. It follows that

\[
o(H-I)=b=|I|+f.
\]

Since \(2|M|=k-f\), this proves the equality in Theorem 1.

For completeness, the equality is a self-contained optimality certificate. For any proposed set \(A\), each odd component of \(H-A\) needs either an unmatched vertex or a matching edge to \(A\). At most \(|A|\) components can use such matching edges. Therefore every matching has size at most \((k-o(H-A)+|A|)/2\). A feasible matching attaining that bound is maximum. The independent verifier uses only this argument and the separately proved exact kernel conditions.

\hypertarget{af-43f6530146-why-the-outer-set-is-exactly-the-exposable-set}{%
\subsection{5. Why the outer set is exactly the exposable set}\label{af-43f6530146-why-the-outer-set-is-exactly-the-exposable-set}}

Section 3 shows that every vertex of \(O\) can be exposed by a maximum matching. Conversely, any matching leaves at least \(b-|I|=f\) vertices of \(O\) unmatched: each odd outer block needs an unmatched vertex or a distinct matching endpoint in \(I\). A maximum matching has exactly \(f\) exposed vertices in the whole graph, so it cannot expose a vertex outside \(O\). Thus \(O\) is exactly the set of vertices exposable by a maximum matching.

A vertex is exposable by a maximum matching precisely when \(\nu(H-v)=\nu(H)\). Equivalently it is the endpoint of a simple even \(M\)-alternating path from an exposed root. One implication follows by flipping that path; the other follows by taking the symmetric difference of \(M\) with a maximum matching that exposes the vertex. Both matchings are maximum, so their path components cannot augment either matching. This proves the remaining claims of Theorem 1 without invoking an external decomposition theorem.

\hypertarget{af-43f6530146-cost-and-executable-evidence}{%
\subsection{6. Cost and executable evidence}\label{af-43f6530146-cost-and-executable-evidence}}

Every vertex becomes outer at most once and has its adjacency scanned once. A matched edge can create at most one inner singleton during the whole search. When its first endpoint becomes inner, the other endpoint becomes outer, and outer status is permanent. If that inner endpoint is later absorbed by a blossom, both endpoints remain outer; the edge cannot create a second inner vertex. Thus there are at most \(\nu=|M|\) inner-vertex births.

Every nontrivial contraction has at least one inner singleton on its odd quotient cycle. All such vertices become outer permanently. Charge each contraction to one of these absorbed inner vertices. No vertex can be charged again, so there are at most \(\nu\) contractions. Ancestor walks, branch rewrites and the full-array contraction scan each cost \(O(k)\). Initialization and ordinary adjacency scans cost \(O(k+h)\). The total is \(O(h+k(\nu+1))\), with linear additional storage. Constructing \(I\) and checking its odd-component count also takes \(O(k+h)\) work.

For the \href{https://github.com/icekylinx/algorithm-frontier/blob/main/problems/general-matching/reference/cover-compression-proof.md}{two-stage reference solver}, the first edge kernel has \(h,k=O(c^2)\) and \(\nu\leq c\). Extraction on that original kernel therefore costs \(O(c^3)\), even when the blossom solve took place on its smaller \(2c\)-vertex induced graph.

The \href{https://github.com/icekylinx/algorithm-frontier/blob/main/problems/general-matching/verify/check_certificate_search.py}{dedicated search checker} compares the actual production function against the deletion-defined exposable set using independent subset dynamic programming. It checks all simple graphs through six vertices and 500 seeded graphs on seven through thirteen vertices. It also verifies the terminal count, the factor-criticality of each outer block, its unique external matched edge, the perfect matching on \(U\), and the exact odd-component count. For graphs through nine vertices it independently enumerates simple even alternating paths. Rejection of a nonmaximum input matching uses explicit production exceptions and remains active under optimized Python. The \href{https://github.com/icekylinx/algorithm-frontier/blob/main/problems/general-matching/verify/check_cover_compression.py}{compression checker} additionally checks the \(\nu\) contraction bound on the production search after lifting the smaller-graph matching back to the original kernel.

The \href{https://github.com/icekylinx/algorithm-frontier/blob/main/problems/general-matching/verify/check_reference.py}{full reference test suite} additionally checks every public solver result with the independent certificate verifier, including native loops, parallel records, isolated vertices, invalid certificates and CLI execution. These are finite implementation checks supporting the proof and review; they are not proof-assistant formalization.
